\documentclass[11pt,a4paper]{article}

\usepackage[utf8]{inputenc}
\usepackage[main=italian,provide=*]{babel}
\usepackage{amsmath}
\usepackage{amssymb}
\usepackage{mathtools}
\usepackage{geometry}
\usepackage{booktabs}
\usepackage{enumitem}
\usepackage{verbatim}
\usepackage{hyperref}
\hypersetup{colorlinks=true,linkcolor=black,citecolor=black,urlcolor=blue}

\geometry{margin=2.7cm}

\title{Perché L-BFGS-B\\
\large Giustificazione numerica e riferimento personale (non parte del testo della tesi)}
\author{}
\date{}

\begin{document}

\maketitle

\noindent\textbf{Nota.} Questo documento non va nella tesi: è il mio riferimento di teoria per una scelta implementativa (il secondo stadio di ottimizzazione nel VQE del trimero ad anello) che nella relazione finale riporto solo nei suoi risultati. Lo tengo qui per ricostruire, quando serve, perché quella scelta è giustificata e non solo "perché funziona".

\section{Il problema}

Nel confronto fra famiglie di ansatz su una griglia di $b/J$ (fatto in una fase precedente del progetto), l'ottimizzatore usato è COBYLA con partenze multiple: un metodo \emph{derivative-free}, che non usa il gradiente, e che da solo lascia tipicamente un residuo sull'energia dell'ordine di $10^{-5}$. Per quello scopo — confrontare qualitativamente quale famiglia di ansatz raggiunge quale fedeltà — un residuo di quell'ordine è ininfluente: non cambia la classifica fra ansatz.

Il punto di lavoro del VQE del trimero ad anello ($J=1$, $J'=0.4$, $b=b_c=2.4$, $D=0.15$, Opzione B) è diverso: lo stato ottimizzato diventa l'ingresso del circuito che calcola le correlazioni dinamiche $C_{ij}^{\alpha\beta}(t)$. La domanda corretta non è quanto vale il residuo sull'\emph{energia}, ma quanto si propaga sull'\emph{osservabile che mi interessa}, cioè il correlatore. Le due cose non sono la stessa cosa, e la sezione successiva quantifica perché.

\section{Dal residuo di energia all'errore sul correlatore}

\subsection{Setup spettrale}

Scrivo lo stato variazionale ottimizzato nella base degli autostati esatti $\{|n\rangle\}$ dell'Hamiltoniana, con energie $E_0<E_1\le E_2\le\dots$:
\[
|\psi(\theta)\rangle = c_0|0\rangle + \sum_{n\ge1} c_n|n\rangle, \qquad \sum_n|c_n|^2=1.
\]
La fidelity rispetto al vero stato fondamentale è $F=|c_0|^2$.

\subsection{Un bound esatto: $1-F\le \Delta E/\Delta$}

L'energia di aspettazione si scrive
\[
E(\theta)-E_0 = \sum_{n\ge1}|c_n|^2(E_n-E_0).
\]
Ogni termine della somma è pesato per almeno il gap al primo eccitato $\Delta \equiv E_1-E_0$, quindi
\[
E(\theta)-E_0 \;\ge\; \Delta\sum_{n\ge1}|c_n|^2 \;=\; \Delta\,(1-F).
\]
Definendo $\Delta E \equiv E(\theta)-E_0$ (il residuo di energia rispetto al fondamentale esatto), questo dà
\[
\boxed{1-F \;\le\; \frac{\Delta E}{\Delta}.}
\]
È un bound esatto, non una stima euristica: segue solo dalla decomposizione spettrale e dal fatto che $E_n-E_0\ge\Delta$ per ogni $n\ge1$.

\subsection{Perché l'energia non è la misura giusta dell'errore}

Definisco $\varepsilon\equiv\sqrt{1-F}$, l'ampiezza "persa" sugli stati eccitati. Dal bound sopra, $\Delta E \ge \Delta\,\varepsilon^2$: il residuo di energia è quadratico in $\varepsilon$. Questo è un fatto generale, non specifico di questo problema: l'energia è stazionaria rispetto a variazioni dello stato attorno al minimo (è proprio la definizione di minimo variazionale), quindi la sua sensibilità a un errore di ampiezza $\varepsilon$ è del second'ordine, $O(\varepsilon^2)$.

Un osservabile generico $O$ (come un correlatore $C_{ij}^{\alpha\beta}$) non ha questa proprietà. Sviluppando:
\[
\langle O\rangle_\psi - \langle O\rangle_0 = -\varepsilon^2 O_{00} + 2\,\mathrm{Re}\Big[c_0^*\sum_{n\ge1}c_n O_{0n}\Big] + O(\varepsilon^2),
\]
dove $O_{mn}=\langle m|O|n\rangle$. Il termine centrale è lineare in $\varepsilon$ (i coefficienti $c_n$ per $n\ge1$ sono essi stessi $O(\varepsilon)$) e, a meno di una simmetria che annulli tutti gli elementi di matrice $O_{0n}$ verso lo stato fondamentale, è generico e non si cancella. Quindi:
\[
\text{errore su } E \sim O(\varepsilon^2)=\Delta E, \qquad\qquad \text{errore su } \langle O\rangle \sim O(\varepsilon) = O(\sqrt{\Delta E/\Delta}).
\]
Un residuo "piccolo" sull'energia produce un errore sull'osservabile che scala come la sua \emph{radice quadrata} — molto più grande in valore assoluto.

\section{Applicazione numerica al punto di lavoro}

Al punto di lavoro del trimero ad anello ($J=1$, $J'=0.4$, $b=b_c=2.4$, $D=0.15$, Opzione B) il gap vero è stato misurato direttamente (diagonalizzazione esatta, non una stima):
\[
\Delta = E_1-E_0 = 0.2547.
\]
Con un residuo COBYLA tipico $\Delta E\approx10^{-5}$ (l'ordine di grandezza riscontrato nelle altre scansioni del progetto quando COBYLA lavora da solo, senza raffinamento successivo):
\begin{align}
1-F &\le \frac{\Delta E}{\Delta} \approx \frac{10^{-5}}{0.2547} \approx 4\times10^{-5},\\
\varepsilon &= \sqrt{1-F} \approx \sqrt{4\times10^{-5}} \approx 6\times10^{-3}.
\end{align}
Quindi l'errore atteso su un correlatore generico calcolato a partire da uno stato ottimizzato solo con COBYLA è dell'ordine di $6\times10^{-3}$ — tre ordini di grandezza più grande del residuo di partenza sull'energia ($10^{-5}$).

Per dargli un metro di paragone concreto: nella relazione delle correlazioni dinamiche, l'autocorrelazione $C_{33}^{zz}(t)$ allo stesso punto di lavoro ha un'oscillazione reale di ampiezza $\approx0.0015$ attorno al suo valore medio — "quasi congelata" nel caso ideale. Un errore sistematico di $6\times10^{-3}$, circa quattro volte più grande di quell'oscillazione, l'avrebbe completamente coperta: non avrei potuto distinguere il segnale fisico dal rumore introdotto dall'ottimizzatore.

\begin{center}
\renewcommand{\arraystretch}{1.3}
\begin{tabular}{@{}lc@{}}
\toprule
\textbf{Quantità} & \textbf{Valore}\\
\midrule
gap $\Delta=E_1-E_0$ (misurato) & $0.2547$\\
residuo tipico COBYLA-solo $\Delta E$ & $\sim10^{-5}$\\
bound sulla fidelity $1-F\le\Delta E/\Delta$ & $\lesssim4\times10^{-5}$\\
ampiezza persa $\varepsilon=\sqrt{1-F}$ & $\approx6\times10^{-3}$\\
oscillazione reale di $C_{33}^{zz}(t)$ (riferimento) & $\approx0.0015$\\
\bottomrule
\end{tabular}
\end{center}

\section{La soluzione: L-BFGS-B}

L-BFGS-B (\emph{Limited-memory Broyden--Fletcher--Goldfarb--Shanno with Bound constraints}) è un metodo quasi-Newton a memoria limitata, con vincoli di box sui parametri. A differenza di COBYLA, usa il gradiente: qui è calcolabile in modo esatto, senza rumore statistico, perché lavoro su simulatore statevector e non su shot reali. Vicino a un minimo, un metodo quasi-Newton converge in modo superlineare, mentre COBYLA (che costruisce solo un modello lineare locale della funzione obiettivo) converge più lentamente. È la combinazione naturale: COBYLA esplora globalmente senza bisogno del gradiente, L-BFGS-B rifinisce localmente sfruttandolo.

\subsection{Riferimenti}

\begin{itemize}[leftmargin=1.3em]
\item R. H. Byrd, P. Lu, J. Nocedal, C. Zhu, \emph{A Limited Memory Algorithm for Bound Constrained Optimization}, SIAM J. Sci. Comput. \textbf{16}(5), 1190--1208 (1995). Non open access: \url{https://epubs.siam.org/doi/10.1137/0916069}. Esiste una versione tecnica preliminare (Argonne National Laboratory / Northwestern University, 1994, tre soli autori — senza Zhu, entrato come coautore solo nella versione pubblicata), liberamente leggibile: \url{https://scispace.com/pdf/a-limited-memory-algorithm-for-bound-constrained-3pm6n4lpp1.pdf}.
\item C. Zhu, R. H. Byrd, P. Lu, J. Nocedal, \emph{Algorithm 778: L-BFGS-B: Fortran subroutines for large-scale bound-constrained optimization}, ACM Trans. Math. Softw. \textbf{23}(4), 550--560 (1997) — le subroutine Fortran usate (indirettamente, via scipy). Non open access: \url{https://dl.acm.org/doi/10.1145/279232.279236}.
\item Documentazione scipy (liberamente accessibile, descrive l'implementazione effettivamente chiamata dal codice): \url{https://docs.scipy.org/doc/scipy/reference/optimize.minimize-lbfgsb.html}.
\end{itemize}

\section{Implementazione}

Codice realmente usato nel progetto (schema a due stadi, identico nella struttura sia per il dimero che per il trimero ad anello):

\begin{verbatim}
def vqe_multistart(ansatz, H_op, w, v, R=6, maxiter=300, seed=0):
    rng = np.random.default_rng(seed)
    n = ansatz.num_parameters
    finals = []
    best = (np.inf, None, None)

    def obj(x):
        sv = Statevector(ansatz.assign_parameters(x))
        return float(np.real(sv.expectation_value(H_op)))

    for _ in range(R):
        x0 = rng.uniform(0.0, 2 * np.pi, n)
        res = minimize(obj, x0, method="COBYLA", options={"maxiter": maxiter})
        finals.append(res.fun)
        if res.fun < best[0]:
            best = (res.fun, res.x, res.nfev)
    Ebest, xbest, nfev = best

    # polish: L-BFGS-B dal miglior punto COBYLA
    res_polish = minimize(obj, xbest, method="L-BFGS-B")
    if res_polish.fun < Ebest:
        Ebest, xbest = res_polish.fun, res_polish.x

    sv = Statevector(ansatz.assign_parameters(xbest))
    fid = ground_fidelity(sv.data, w, v)
    return {"E": Ebest, "x": xbest, "fid": fid, "nfev": nfev,
            "finals": np.array(finals), "spread": np.ptp(finals),
            "statevector": sv}
\end{verbatim}

\section{Nota sull'uso di un assistente AI}

La chiamata a scipy che implementa il secondo stadio (\texttt{scipy.optimize.minimize} con \texttt{method="L-BFGS-B"}) l'ho fatta scrivere a un assistente AI, dopo avergli chiesto se esistesse un metodo adatto allo scopo — portare a convergenza fine un punto già individuato da un ottimizzatore globale, sfruttando che qui il gradiente è disponibile esattamente. Il conto delle Sezioni 2--3 — perché un residuo "piccolo" sull'energia non lo è affatto su un correlatore, e di quanto — l'ho fatto io, prima ancora di cercare la soluzione: è quello che mi ha convinto che il problema fosse reale e non solo un timore teorico.

\end{document}
